% !PRINTED 461
\weberpage{461}

If we substitute for $t$ in $\chi(t)$ the various quantities (4), we
obtain $n$ quantities, one from each of the conjugate fields
\weq{6}{\Theta = \chi(\alpha),\ \Theta_1 = \chi(\alpha_1) \ldots \Theta_{n-1} = \chi(\alpha_{n-1}),}
and these are called the \gesperrt{quantities conjugate to} $\Theta$.

Every symmetric function of these quantities is at the same time a
symmetric function of the roots of the equation (3), and hence
contained in $\Omega$.

Among these symmetric functions we wish to single out two, which
occur frequently, by special names and notations; there is the sum
\weq{7}{S(\Theta) = \Theta + \Theta_1 + \Theta_2 + \cdots + \Theta_{n-1},}
which we call the \gesperrt{trace} of $\Theta$, and the product
\weq{8}{N(\Theta) = \Theta \Theta_1 \Theta_2 \ldots \Theta_{n-1},}
which is called the \gesperrt{norm} of $\Theta$. Conjugate numbers accordingly
have the same traces and norms.

The product
\weq{9}{(t - \Theta)(t - \Theta_1) \ldots (t - \Theta_{n-1}) = \Phi(t)}
is an integral rational function of degree $n$ of $t$ in
$\Omega$, and its roots are $\Theta$ and the quantities conjugate to
$\Theta$. From this results the theorem:

\gesperrt{1. Every quantity in $\Omega(\alpha)$ is a root of an
equation of degree $n$ in $\Omega$, whose remaining roots
are the quantities conjugate to $\Theta$.}

The computation of the function $\Phi(t)$ from $f(x)$ is nothing
other than the \gesperrt{Tschirnhausen transformation} considered
from another point of view in the fourth and sixth sections.

We now have to examine the function $\Phi(t)$ as to its
irreducibility.

If the function $\Phi(t)$ is reducible, then it has an
irreducible factor $\varphi(t)$, in which we may assume the coefficient of the
highest power of $t$ equal to 1, and every such
factor vanishes for at least one of the values $\Theta,
\Theta_1 \ldots \Theta_{n-1}$.

Let us then have
\[
\varphi(\Theta) = \varphi[\chi(\alpha)] = 0,
\]
that is, the equations $\varphi[\chi(x)] = 0$ and $f(x) = 0$ have
a common root. But since $f(x)$ is assumed irreducible,
$\varphi[\chi(x)]$ must, by the Theorem II, \S.\,141, for all roots
%CONT>>
